\paragraph{The maximal slope of the full cusp lattice.}
\label{par:parent-full-cusp-slope}
The full cusp lattice and its row have different roles in this field comparison. Let \(K_c\)
be the field of the geometric cusp model at the base level,
and \(K\) the field carrying the row. The size estimate for the
recursive class sum bounds \(\log\operatorname{rd}(K_c)\) by
\(217(frL)^2\), since the field of the core is the coefficient
field in that estimate. We first work over \(K_c\), which is
totally complex. Write \(\|\cdot\|_{\rm Pet}\) for the adelic probability
Petersson norm multiplied by the Lie-frame factor \(a_i\)
above. Use the same factor on sup norms. It cancels in their
comparison. The integral lattice
has full rank there. A nonzero good component is supplied by the
base realization of rank \(D_0\). Thus every application of a
maximal-slope or a logarithmic rank bound below has positive
ambient rank. \label{cor:parent-cusp-slope-1}

\begin{established}[Sizes on the completed parent]\label{hyp:parent-cusp-size}
Use the geometric field \(K_c\), row field \(K\), and normalized
Petersson metric in the common primitive frame. Choose a fixed exponent \(a\) with \(frn_V L\le X=x^a\le b\) on the late family, where
\(L=\log b\ge10\), \(f\ge8\) and \(r\ge2\).
Write \(s\) for the sum of the absolute frame powers, \(t=frn_V\)
and \(w=1+s+t\). Summing the dominant weight bounds on the
two signature blocks gives \(s\le8frb\), hence
\(w\le10frb\) and \(\log w\le4L\).
With \(\ell=2L\), the two-prime level \(M=g_0g_1\), the
quadratic discriminant parameter \(d\) and \(d_E=|\operatorname{disc}E|\)
satisfy
\[
 \log(fr)\le\ell,\quad \log M\le30\ell,\quad
 \log d_E\le4f\ell,\quad \log d\le4\ell,\quad \log w\le2\ell.
\]
The marked CM period satisfies \(|\log\Omega|\le3L\).
\end{established}
\begin{proof}
The ranks of Lemma~\ref{lem:window-parameter-choice} and the number of columns in the completed point grow at most polynomially in \(x\); \(b=\exp(x^\epsilon)\) eventually dominates each fixed power of \(x\). The upper box bound \(2b\) per entry, applied to the two frame lists and their determinant shifts, gives \(s\le8frb\). On the selected family \(frn_VL\le x^a\le b\), hence \(w\le10frb\) and \(\log w\le4L\) for large \(x\). For the same field and two primes the calculation of Lemma~\ref{lem:deep-cover-field-size} gives \(\log g_j\le9L+1\), and its bound for the root discriminant gives \(\log |\operatorname{disc} E|=2f\log\operatorname{rd}(E)\le 4f\ell\). The two-prime product and the marked discriminant then give the remaining logarithmic bounds. The period bound is the CM period estimate on the retained family used in Lemma~\ref{lem:terminal-centred-weights}.
\end{proof}

\begin{proposition}[A uniform upper bound for the parent maximal slope]
\label{prop:parent-cusp-maximal-slope}
For each lift of the completed parent with
Lemma~\ref{hyp:parent-cusp-size}, the ambient integral cusp
lattice satisfies
\[
  \mu_{\max}(\mathrm{lat}_{\rm all})\le x^A b
\]
for a fixed \(A\), uniformly in the row field. The maximal slope
uses the normalized Petersson metric and the actual ambient
lattice in the primitive frame.
\end{proposition}

\begin{proof}
\proofstep{The lower norm bound for an integral cusp form.}
For a nonzero integral form, apply the CM reading
\eqref{eq:slice-cm-point-norm} with these powers. The cusp
flow at level \(M\) uses the thickness
\(e_0=\min(1,\pi/(2M\sqrt d(fw+2)\log2))\).
The CM approximation there gives a frame bound
\(\Lambda=R^{2q+1}\), where \(q\le r\) is the negative signature
on this lift. To use the negative block, approximate its \(q\)
axes in the signed Hermitian space and take the orthogonal
complement. The graph bound is invariant under passage to its
adjoint. The Gram determinant which clears the complementary
projection then has size at most \((2A)^{2q}\) in the notation
of the CM reading. This gives
\[
 R=\tfrac{1600}{3}r\sqrt{fw}\,(fr)^3d_E/\sqrt{e_0}.
\]
Indeed \(-\log e_0\le\log(2M\sqrt d f w)\).

Taking logarithms in the formula for \(R\) gives
\(\log R\le(4f+32)\ell\). The nonzero integral CM
coefficient, with its dual letter metrics, now gives
\(\operatorname{mean}_\sigma\log\|F\|_{\sup,\sigma}\ge-\lambda_s\),
where
\[
 \lambda_s\le
 \bigl((1+|\log\Omega|)/2+\ell+
                  (2q+1)(6f+48)\ell\bigr)(s+t).
\]
Here the value line contributes half the logarithm of
\(\Lambda\) per letter, and reading the frame coefficient costs
at most \(s\log(r\Lambda)\). \label{cor:parent-cusp-slope-2}

\proofstep{Returning to the Petersson norm.}
For the passage back to the Petersson norm, let \(B\) be the
bound on squared point evaluation. In a flag chart of complex
dimension \(D=fr(r-1)/2\) take \(h=(200r^3(w+f))^{-1}\).
On the part with larger weights the
holomorphic submean inequality contributes
\(8D!/(\pi h^2)^D\), including the variation of the frame norm
and the choice of a point above half the supremum. Overlap
of the charts is bounded by the number of integral transition
matrices moving the frame on the subspace with larger weights into a fixed ball. Packing their
entries as in the evaluation-net argument bounds this
multiplicity by
\[
 \mathcal M(e_0)=(1+2(2fr)^2d_E/e_0)^{2fr^2}.
\]
Integration on the
probability space introduces the upper level mass \(V_M\)
from the measure comparison. Thus one can take
\[
 B=\frac{8D!}{(\pi h^2)^D}\mathcal M(e_0)V_M,\qquad
             \|F\|_{\sup,\sigma}^2\le B\|F\|_{\rm Pet,\sigma}^2 .
\]

For clarity, the chart factor has logarithm at most
\(3\log2+fr^2(12\ell+16\log2)/2\). In the packing bound use
\(T=3fwM\sqrt d\ge e_0^{-1}\). Its logarithm is at most
\(2\log2+35\ell\). The log multiplicity is at most
\(2fr^2(\log16+2\log(fr)+\log d_E+\log T)\).
The log mass is at most \(130fr^2\ell\).
Adding gives \[
 0\le\log B\le250 f^2r^2\ell.
\]
These bounds use the equivariant frame character at the level,
the even central weight, and the same Haar probability as the norm.
\label{cor:parent-cusp-slope-3}

\proofstep{The rank bound using the same finite net.}
The dimension bound uses this identical \(B\). At an embedding
of \(K_c\), a rationally independent frame of forms stays
independent in the space of holomorphic readings by the
descent in paragraph~\ref{par:slice-evaluation-net}. If \(F_1,\ldots,F_d\) is an orthonormal basis for the
Petersson probability measure \(\mu\), the evaluation functional
at \(z\) has squared norm \(\sum_i\|F_i(z)\|^2\). The evaluation
inequality applies to every vector in their span, so
\[
 d=\int\sum_{i=1}^d\|F_i(z)\|^2\,d\mu(z)
       \le\int B\,d\mu(z)=B.
\]
This is the same \(B\) as in the evaluation bound.

Now in any nonzero integral sublattice of rank \(u\) and
squared-norm slope \(\mu\), the Minkowski bound used for the
short section gives a nonzero form with
\[
 2\operatorname{mean}_{\sigma}\log\|F\|_{\rm Pet,\sigma}
      \le -\mu+\log u+\log\operatorname{rd}(K_c).
\]
Indeed the Euclidean ball of radius one in each complex rank
\(u\) contains the product of coordinate discs of radii
\(1/\sqrt u\), and the covolume of the coefficient integers
contributes the square root of the discriminant per rank.
Taking the two averages with the same place multiplicities
gives
\[
 \mu\le 2\lambda_s+2\log B+\log\operatorname{rd}(K_c)
              \le x^{2a+4}b .
\]
For the last estimate \(2q+1\le3r\), \(6f+48\le12f\) and
\(frL\le X\) give
\[
 \lambda_s\le760X^2 b.
\]

The other two contributions together are bounded by
\[
 1000X^2+217X^2=1217X^2.
\]
The factor \(x^4\) absorbs their sum
and the factor two on \(\lambda_s\) for \(x\ge10\).
\label{cor:parent-cusp-slope-4}

\proofstep{The actual ambient arithmetic lattice.}
On the actual row the ambient integral forms are the
\(\mathcal O_K\)-span of this lattice under the independent
semilinear realization. The norm of every geometric vector
at an extension \(\sigma\) is its geometric norm multiplied
by the positive scalar \(c_\sigma\) of the measure comparison
above. The direct and inverse comparisons amount to
\[
 e^{-122U}\le c_\sigma\le e^{122U}.
\]

Divide out this scalar. The normalized Gram on geometric
vectors extends to their complex span, by polarization and
density of the image of the totally complex field. Choosing
the conjugate embedding conjugates the Gram and preserves
its determinant. For the bound on arbitrary sublattices, let \(K'/K_c\) be a
normal extension containing \(K\), and write \(M_c\) for the
geometric lattice over \(K_c\). Give \(M_c\otimes\mathcal O_{K'}\)
the normalized metric. The Galois group permutes its places
by isometries, hence preserves the maximal destabilizing
saturated lattice. Its rational span descends to \(V_c\) over
\(K_c\).

To identify its integral lattice, tensor the exact
kernel sequence
\[
 0\longrightarrow M_c\cap V_c\longrightarrow M_c
       \longrightarrow (M_c\otimes K_c)/V_c
\]
with \(\mathcal O_{K'}\) over \(\mathcal O_{K_c}\).
Flatness identifies its kernel with the saturated
intersection over \(K'\). Normalized slopes of these extended
lattices are unchanged by the ramification and archimedean
multiplicity identities. Any sublattice over \(K\) also extends
at its original normalized slope. The maximal destabilizing
lattice over \(K'\) therefore bounds the slope over \(K\).
Restoring \(c_\sigma\) changes any squared slope by at most
\(244U\). Since \(U=fmbL\log\log x\le x^{a+1}b\) on our
family by \(m\le rn_V\) and \(\log\log x\le x\), we have
\(\mu_{\max}(\mathrm{lat}_{\rm all})\le x^A b\) for a fixed
\(A\), uniformly in the row field and the lift.
\label{cor:parent-cusp-slope-5}

\end{proof}

\paragraph{Rank on the two sides of the projection.}
The metric comparison leaves a rank factor on each side. To identify these ranks, use
the central projection to the retained tuples of
semisimple labels in the same finite group at the base level.
In the rational form frame the type space is the intersection
of the kernels of the character equations. Adding the kernel
of the retained projector gives the discarded space. Flatness
and the fact that the full integral forms span compute the two
lattice ranks from exactly these linear equations after any
complex embedding. Twisting by the inverse central form
character makes the first system the Whittaker character
equations. The group order clears the projection on the full
lattice, since its character coefficients times the order are
algebraic integers. Inversion conjugates these coefficients at
each embedding and the group acts unitarily, so the two parts
are orthogonal. \label{cor:parent-cusp-slope-6}

\begin{lemma}[Rank discarded in the primitive base frame]
\label{lem:parent-projector-rank-loss}
Use the retained central projector in
Proposition~\ref{prop:equivariant-selected-row}, transported to the
primitive base frame, with the finite label bounds of
Proposition~\ref{prop:paired-good-rank}. Put \(\eta=\exp(-x^{e_z/7})\).
The retained and discarded ranks \(d_g,d_b\) in the actual type lattice satisfy
\[
 d_b\le3\eta(d_b+d_g),\qquad d_g=D_0.
\]

\end{lemma}

\begin{proof}
\proofstep{Computing the finite bad-rank fraction.}
In the definite model the two groups \(H_k\) in
\eqref{eq:row-bad-share} act freely on the marked components.
On a nonterminal deep cover use the inverse images of \(P\)
(the character subgroup) and of \(H_k\) in its level
automorphism group. The pullback action on forms differs
from translation by a single character. Thus the product of
a translated separating section and a translated form still
has this character on the form factor. The invariant-sum map
in the proof of \eqref{eq:row-bad-share} embeds the restriction
to each inverse image \(H\) in \(a_H\) regular representations.

Let \(r_+,r_-\) be the ranks at weights \(\kappa+2c_0\) and \(\kappa-c_0\). The determinant injections and the rank growth at shift two on this cover give
\[
 a_H|H|\le r_+\le\tfrac32 r_-\le\tfrac32|\pi^{-1}P|\,d_{\rm type}.
\]
Here \(c_0\) is the weight
of the orbit-separating section on that cover, with the group
determinant nonzero on every component. The regular trace on
the inverse image reads the local semisimple-label share
downstairs, since the averaging factor for the kernel cancels
in the trace ratio. Using \eqref{eq:row-bad-types} for both
primes, each discarded share is at most
\(\eta=\exp(-x^{e_z/7})\) on the late family. Hence
\[
   d_b\le (3/2)(\eta+\eta)d_{\rm type}
       =3\eta(d_b+d_g).
\]
The definite regular comparison has factor one and satisfies
the same inequality. The retained image is the good space in
the primitive base frame, of rank \(d_g=D_0\).
\label{cor:parent-cusp-slope-7}
\end{proof}
